\documentclass[UTF8,a4paper,11pt]{ctexart}
\usepackage[a4paper,margin=2.2cm]{geometry}
\usepackage{amsmath,amssymb,amsthm}
\usepackage{booktabs,longtable,array}
\usepackage{hyperref}
\usepackage{xcolor}
\usepackage{listings}
\usepackage{enumitem}
\usepackage{float}
\usepackage{multirow}
\usepackage{caption}

\hypersetup{
  colorlinks=true,
  linkcolor=blue!60!black,
  urlcolor=blue!60!black
}

\lstset{
  language=Python,
  basicstyle=\ttfamily\small,
  keywordstyle=\color{blue!70!black},
  commentstyle=\color{green!40!black},
  stringstyle=\color{orange!50!black},
  showstringspaces=false,
  breaklines=true,
  frame=single,
  backgroundcolor=\color{gray!5},
  columns=fullflexible
}

\setcounter{tocdepth}{3}

\title{CS336 Assignment 1 实现报告}
\author{项目路径：\texttt{assignment1-basics}}
\date{\today}

\begin{document}
\maketitle
\tableofcontents
\newpage

\section{Assignment 1 概述}

\subsection{完成内容、最终目的与总体结论}
本次 Assignment 1 的核心目标，是从零搭建一个\textbf{decoder-only Transformer LM} 的最小可运行实现链条：先用 \textbf{byte-level BPE} 完成文本到 token 的映射，再实现训练阶段所需的\textbf{数据加载、损失函数、优化器、学习率调度、梯度裁剪与 checkpointing}，最后实现 \textbf{注意力、RoPE、Transformer Block 与 Transformer LM}。从工程结构上看，这个作业实际上把一个小型 LLM 训练系统拆成了若干可独立验证的组件，并要求每个组件都能在统一的测试入口上工作。

本报告聚焦于 \textbf{代码型 deliverables}。也就是说，本报告覆盖了下列 handout 中与代码直接相关的部分：
\begin{itemize}[leftmargin=2em]
  \item \texttt{train\_bpe}, \texttt{tokenizer}
  \item \texttt{linear}, \texttt{embedding}, \texttt{rmsnorm}, \texttt{positionwise\_feedforward}
  \item \texttt{rope}, \texttt{softmax}, \texttt{scaled\_dot\_product\_attention}, \texttt{multihead\_self\_attention}
  \item \texttt{transformer\_block}, \texttt{transformer\_lm}
  \item \texttt{cross\_entropy}, \texttt{adamw}, \texttt{learning\_rate\_schedule}, \texttt{gradient\_clipping}
  \item \texttt{data\_loading}, \texttt{checkpointing}
\end{itemize}

最终结果是：本地完整测试集达到
\[
\boxed{47\ \text{passed},\ 1\ \text{xfailed}}
\]
其中唯一的 \texttt{xfailed} 是测试中\textbf{显式标记的预期失败项}，并不表示实现错误。完整测试命令的 wall-clock 时间约为 \(15.04\) 秒，而 pytest 报告的测试执行时间为 \(13.04\) 秒。

\subsection{关键测试命令}
表 \ref{tab:key-commands} 给出了本次实现过程中最常用的验证命令。

\begin{table}[H]
\centering
\caption{关键测试命令}
\label{tab:key-commands}
\begin{tabular}{p{0.34\textwidth} p{0.58\textwidth}}
\toprule
命令 & 用途 \\
\midrule
\texttt{./.venv/bin/python -m pytest -q} & 运行完整测试集，验证所有组件集成后的最终状态。 \\
\texttt{./.venv/bin/python -m pytest -k test\_train\_bpe\_speed} & 专门验证 BPE 训练的速度约束。 \\
\texttt{./.venv/bin/python -m pytest -k test\_get\_lr\_cosine\_schedule} & 验证 warmup + cosine decay 学习率调度。 \\
\texttt{./.venv/bin/python -m pytest -k 'test\_rope or test\_multihead\_self\_attention\_with\_rope'} & 验证 RoPE 与带 RoPE 的 MHA。 \\
\texttt{./.venv/bin/python -m pytest -k 'test\_transformer\_block or test\_transformer\_lm'} & 验证 Block 与 LM 级别的集成正确性。 \\
\texttt{./.venv/bin/python -m py\_compile ...} & 在运行单测前先做语法级静态检查。 \\
\bottomrule
\end{tabular}
\end{table}

\subsection{代码结构概览}
最终代码主要位于 \texttt{cs336\_basics} 目录中，形成了如下功能分层：
\begin{itemize}[leftmargin=2em]
  \item \textbf{Tokenizer 层}：\texttt{tokenizer.py}
  \item \textbf{基础张量算子}：\texttt{linear.py}, \texttt{embedding.py}, \texttt{silu.py}, \texttt{rmsnorm.py}, \texttt{swiglu.py}
  \item \textbf{注意力层}：\texttt{softmax.py}, \texttt{scaled\_dot\_product\_attention.py}, \texttt{rope.py}, \texttt{multihead\_self\_attention.py}, \texttt{multihead\_self\_attention\_with\_rope.py}
  \item \textbf{模型组装层}：\texttt{transformer\_block.py}, \texttt{transformer\_lm.py}
  \item \textbf{训练工具层}：\texttt{get\_batch.py}, \texttt{cross\_entropy.py}, \texttt{gradient\_clipping.py}, \texttt{AdamW.py}, \texttt{lr\_schedule.py}, \texttt{save\_load\_checkpoint.py}
\end{itemize}

\newpage
\section{组件级实现说明}

\subsection{Tokenizer 与数据预处理}

\subsubsection{组件：\texttt{train\_bpe}}
\paragraph{1. 作用、输入与输出}
该组件用于从原始文本语料训练 \textbf{byte-level BPE}。输入为文本路径 \texttt{input\_path}、目标词表大小 \texttt{vocab\_size} 与可选的 \texttt{special\_tokens}；输出为
\[
\texttt{vocab}: \{ \texttt{id} \mapsto \texttt{bytes} \}, \qquad
\texttt{merges}: [(b_1,b_2), \dots]
\]
其中 \texttt{vocab} 是 token id 到字节串的映射，\texttt{merges} 是按学习顺序记录的 merge 规则。

\paragraph{2. 代码实现摘要}
核心代码如下：
\begin{lstlisting}
token_counts = _collect_pretoken_counts(text, normalized_special_tokens)
pair_counts, pair_to_sequences = _build_pair_statistics(token_counts)

while len(vocab) < vocab_size:
    best_pair = max(pair_counts.items(), key=lambda item: (item[1], item[0]))[0]
    affected_sequences = list(pair_to_sequences.get(best_pair, ()))
    ...
\end{lstlisting}
实现流程是：
\begin{enumerate}[leftmargin=2em]
  \item 用 GPT-2 风格 regex 做 pre-tokenization；
  \item 将每个 pretoken 转为 UTF-8 字节序列；
  \item 统计所有相邻 pair 的频次；
  \item 反复选出最高频 pair 做 merge；
  \item 更新词表与 merge 序列，直到达到目标词表大小。
\end{enumerate}

\paragraph{3. 数学原理}
如果把某个 pretoken 写成 token 序列
\[
t = (u_1, u_2, \dots, u_n),
\]
则 BPE 在每一轮选择
\[
(u^\ast, v^\ast) = \arg\max_{(u,v)} \Big( c(u,v),\ (u,v) \Big),
\]
其中 \(c(u,v)\) 是 pair \((u,v)\) 在整个语料中的总频次；当频次并列时，再按 pair 的字典序做 tie-break。选中后执行
\[
u^\ast, v^\ast \longrightarrow u^\ast v^\ast.
\]
代码中的
\begin{lstlisting}
best_pair = max(pair_counts.items(), key=lambda item: (item[1], item[0]))[0]
\end{lstlisting}
正对应这个规则。

\paragraph{4. 工程技巧}
这一部分用了最重要的工程优化：\textbf{增量维护 pair 统计}。初版做法是在每一轮 merge 后对全部 token sequence 重新统计 pair 频次，复杂度较高，导致 \texttt{test\_train\_bpe\_speed} 失败。最终版本维护了：
\begin{itemize}[leftmargin=2em]
  \item \texttt{pair\_counts}：全局 pair 频次；
  \item \texttt{pair\_to\_sequences}：pair 到包含该 pair 的 sequence 的反向索引。
\end{itemize}
这样每轮只更新受到当前 \texttt{best\_pair} 影响的序列，而不必全量重扫整个语料。

\subsubsection{组件：\texttt{Tokenizer}}
\paragraph{1. 作用、输入与输出}
\texttt{Tokenizer} 负责在\textbf{文本与 token id 序列}之间建立双向映射。输入包括：
\begin{itemize}[leftmargin=2em]
  \item \texttt{vocab: dict[int, bytes]}
  \item \texttt{merges: list[tuple[bytes, bytes]]}
  \item \texttt{special\_tokens: Sequence[str] | None}
\end{itemize}
输出行为则由以下接口提供：
\begin{itemize}[leftmargin=2em]
  \item \texttt{encode(text)}：\(\texttt{str} \to \texttt{list[int]}\)
  \item \texttt{decode(ids)}：\(\texttt{list[int]} \to \texttt{str}\)
  \item \texttt{encode\_iterable(iterable)}：流式编码文本块
  \item \texttt{from\_files(...)}：从 GPT-2 风格词表文件恢复 tokenizer
\end{itemize}

\paragraph{2. 代码实现摘要}
\texttt{encode} 的核心结构是：
\begin{lstlisting}
for is_special, unit in iter_text_units(text, self.special_tokens):
    if is_special:
        token_ids.append(self.special_token_to_id[unit])
    else:
        token_ids.extend(self._encode_pretoken_to_ids(unit))
\end{lstlisting}
\texttt{decode} 的核心结构是：
\begin{lstlisting}
token_bytes = [self.vocab[token_id] for token_id in ids]
return b"".join(token_bytes).decode("utf-8", errors="replace")
\end{lstlisting}
\texttt{encode\_iterable} 则维护一个 \texttt{buffer}，在 chunk 边界上保留不稳定尾部，确保跨 chunk 的 pretoken / special token 不被错误拆开。

\paragraph{3. 数学原理}
编码过程本质上是复合映射：
\[
\text{text}
\xrightarrow{\text{special split}}
\text{units}
\xrightarrow{\text{pretokenize}}
\text{pretokens}
\xrightarrow{\text{BPE merge}}
\text{byte tokens}
\xrightarrow{\text{lookup}}
\text{ids}.
\]
解码过程则是
\[
(i_1, i_2, \dots, i_n)
\xrightarrow{\text{vocab lookup}}
(b_1, b_2, \dots, b_n)
\xrightarrow{\text{concat}}
b
\xrightarrow{\text{utf-8 decode}}
\text{text}.
\]

\paragraph{4. 工程技巧}
这里用了两点工程技巧：
\begin{itemize}[leftmargin=2em]
  \item \textbf{重叠 special token 的长度优先匹配}：通过按长度降序排序 special token，解决短 token 抢先匹配的问题；
  \item \textbf{流式编码的稳定尾部分离}：\texttt{encode\_iterable} 不会简单地对每个 chunk 直接编码，而是保留 chunk 末尾可能跨边界延续的 pretoken。
\end{itemize}

\subsubsection{组件：\texttt{get\_batch}}
\paragraph{1. 作用、输入与输出}
\texttt{get\_batch} 从一维 token 序列中随机抽取训练样本。输入是
\[
\texttt{dataset},\ \texttt{batch\_size},\ \texttt{context\_length},\ \texttt{device},
\]
输出是两个 shape 都为
\[
(\texttt{batch\_size}, \texttt{context\_length})
\]
的 tensor：
\[
X,\ Y.
\]

\paragraph{2. 代码实现摘要}
\begin{lstlisting}
start_indices = np.random.randint(0, num_possible_starts, size=batch_size, dtype=np.int64)
offsets = np.arange(context_length + 1, dtype=np.int64)
batch_np = dataset[start_indices[:, None] + offsets[None, :]]
x = batch[:, :-1]
y = batch[:, 1:]
\end{lstlisting}

\paragraph{3. 数学原理}
设一维 token 流为
\[
(x_0, x_1, \dots, x_{N-1}),
\]
随机抽取起点 \(s_i\) 后，定义：
\[
X_i = (x_{s_i}, x_{s_i+1}, \dots, x_{s_i+L-1}),
\]
\[
Y_i = (x_{s_i+1}, x_{s_i+2}, \dots, x_{s_i+L}),
\]
其中 \(L = \texttt{context\_length}\)。这就是标准 next-token prediction 训练对。

\paragraph{4. 工程技巧}
实现里一次性取出长度 \(L+1\) 的窗口块，然后再切成 \(X\) 与 \(Y\)。这样只做一次 gather；同时 CUDA 路径下使用了 \texttt{pin\_memory} 与 \texttt{non\_blocking=True}。

\subsection{基础算子与前馈网络}

\subsubsection{组件：\texttt{linear}}
\paragraph{1. 作用、输入与输出}
\texttt{linear} 实现线性层。输入是权重
\[
W \in \mathbb{R}^{d_{\text{out}} \times d_{\text{in}}}
\]
和特征
\[
X \in \mathbb{R}^{\cdots \times d_{\text{in}}},
\]
输出
\[
Y \in \mathbb{R}^{\cdots \times d_{\text{out}}}.
\]

\paragraph{2. 代码实现摘要}
\begin{lstlisting}
return in_features @ weights.transpose(-1, -2)
\end{lstlisting}

\paragraph{3. 数学原理}
\[
Y = X W^\top.
\]
因为代码中 \texttt{weights} 的存储形式是 \((d_{\text{out}}, d_{\text{in}})\)，所以在右乘前需要转置。

\subsubsection{组件：\texttt{embedding}}
\paragraph{1. 作用、输入与输出}
\texttt{embedding} 根据 token id 取词向量。输入是
\[
W \in \mathbb{R}^{V \times d_{\text{model}}}, \qquad \texttt{token\_ids} \in \mathbb{N}^{\cdots}
\]
输出
\[
E \in \mathbb{R}^{\cdots \times d_{\text{model}}}.
\]

\paragraph{2. 代码实现摘要}
\begin{lstlisting}
return weights[token_ids]
\end{lstlisting}

\paragraph{3. 数学原理}
这是典型的查表操作：
\[
E_{i} = W[\texttt{token\_id}_i].
\]

\subsubsection{组件：\texttt{silu}}
\paragraph{1. 作用、输入与输出}
逐元素非线性激活，输入输出 shape 完全一致。

\paragraph{2. 代码实现摘要}
\begin{lstlisting}
return in_features * torch.sigmoid(in_features)
\end{lstlisting}

\paragraph{3. 数学原理}
\[
\mathrm{SiLU}(x) = x \cdot \sigma(x), \qquad
\sigma(x)=\frac{1}{1+e^{-x}}.
\]

\subsubsection{组件：\texttt{rmsnorm}}
\paragraph{1. 作用、输入与输出}
对最后一个维度执行 RMSNorm。输入是特征 \(x\) 与可学习缩放参数 \(w\)，输出与输入同 shape。

\paragraph{2. 代码实现摘要}
\begin{lstlisting}
rms = torch.sqrt(torch.mean(in_features * in_features, dim=-1, keepdim=True) + eps)
normalized = in_features / rms
return normalized * weights
\end{lstlisting}

\paragraph{3. 数学原理}
对单个向量 \(x \in \mathbb{R}^{d}\)，RMSNorm 定义为
\[
\mathrm{RMSNorm}(x)=
\frac{x}{\sqrt{\frac{1}{d}\sum_{j=1}^{d}x_j^2+\epsilon}}
\odot w.
\]
代码中的 \texttt{torch.mean(in\_features * in\_features, dim=-1)} 对应分母里的均方，随后再乘以 \texttt{weights} 完成仿射缩放。

\subsubsection{组件：\texttt{swiglu}}
\paragraph{1. 作用、输入与输出}
\texttt{SwiGLU} 是 Transformer FFN 的核心前馈模块。输入是
\[
X \in \mathbb{R}^{\cdots \times d_{\text{model}}},
\]
输出仍然是
\[
\mathbb{R}^{\cdots \times d_{\text{model}}}.
\]

\paragraph{2. 代码实现摘要}
\begin{lstlisting}
w1_out = linear(..., weights=w1_weight, in_features=in_features)
w3_out = linear(..., weights=w3_weight, in_features=in_features)
gated = silu(w1_out) * w3_out
return linear(..., weights=w2_weight, in_features=gated)
\end{lstlisting}

\paragraph{3. 数学原理}
\[
\mathrm{SwiGLU}(x)=W_2\big(\mathrm{SiLU}(W_1x)\odot (W_3x)\big).
\]
其中 \(W_1\) 与 \(W_3\) 都把输入从 \(d_{\text{model}}\) 映射到 \(d_{\text{ff}}\)，而 \(W_2\) 再把中间结果映射回 \(d_{\text{model}}\)。

\subsection{数值计算、优化与训练工程}

\subsubsection{组件：\texttt{softmax}}
\paragraph{1. 作用、输入与输出}
沿指定维度 \texttt{dim} 把任意 logits 转成概率分布，输出 shape 与输入完全一致。

\paragraph{2. 代码实现摘要}
\begin{lstlisting}
max_values = torch.amax(in_features, dim=dim, keepdim=True)
shifted = in_features - max_values
exp_shifted = torch.exp(shifted)
partition = torch.sum(exp_shifted, dim=dim, keepdim=True)
return exp_shifted / partition
\end{lstlisting}

\paragraph{3. 数学原理}
\[
\mathrm{softmax}(x)_i = \frac{e^{x_i}}{\sum_j e^{x_j}}.
\]
为避免数值溢出，代码先做平移：
\[
\tilde{x}_i = x_i - \max_j x_j.
\]
由于 softmax 对所有分量同时减去同一个常数不变，因此
\[
\mathrm{softmax}(x)=\mathrm{softmax}(\tilde{x}).
\]

\paragraph{4. 工程技巧}
这里的关键技巧是\textbf{减去最大值}，这是数值稳定版 softmax 的标准写法。

\subsubsection{组件：\texttt{cross\_entropy}}
\paragraph{1. 作用、输入与输出}
对一个 batch 的 logits 与离散标签计算平均交叉熵。输入：
\[
Z \in \mathbb{R}^{B \times V}, \qquad y \in \{0,\dots,V-1\}^B.
\]
输出是标量 loss。

\paragraph{2. 代码实现摘要}
\begin{lstlisting}
max_values = torch.amax(inputs, dim=-1, keepdim=True)
shifted = inputs - max_values
log_sum_exp = torch.log(torch.sum(torch.exp(shifted), dim=-1)) + max_values.squeeze(-1)
target_logits = inputs.gather(dim=-1, index=targets.unsqueeze(-1)).squeeze(-1)
losses = log_sum_exp - target_logits
return losses.mean()
\end{lstlisting}

\paragraph{3. 数学原理}
单样本交叉熵可写成
\[
\mathcal{L}(z,y) = -\log\frac{e^{z_y}}{\sum_j e^{z_j}}
 = \log\sum_j e^{z_j} - z_y.
\]
代码中的 \texttt{log\_sum\_exp} 对应前一项，\texttt{target\_logits} 对应后一项。

\paragraph{4. 工程技巧}
依然采用了\textbf{log-sum-exp 的稳定写法}，避免直接对大 logits 做指数运算。

\subsubsection{组件：\texttt{gradient\_clipping}}
\paragraph{1. 作用、输入与输出}
对一组参数梯度做 global \(L_2\) norm clipping，输入是一组参数与阈值 \texttt{max\_l2\_norm}，输出为空，但会\textbf{原地修改} \texttt{param.grad}。

\paragraph{2. 代码实现摘要}
\begin{lstlisting}
params_with_grad = [p for p in parameters if p.grad is not None]
total_norm = torch.norm(
    torch.stack([torch.norm(p.grad.detach(), p=2) for p in params_with_grad]),
    p=2,
)
clip_coef = max_l2_norm / (total_norm + 1e-6)
if clip_coef < 1:
    for param in params_with_grad:
        param.grad.mul_(clip_coef)
\end{lstlisting}

\paragraph{3. 数学原理}
设所有梯度为 \(g_1,\dots,g_n\)，则全局范数为
\[
\|g\|_2 = \sqrt{\sum_{i=1}^{n}\|g_i\|_2^2}.
\]
若 \(\|g\|_2 > \tau\)，则统一缩放：
\[
g_i \leftarrow g_i \cdot \frac{\tau}{\|g\|_2}.
\]

\paragraph{4. 工程技巧}
使用 \texttt{mul\_} 做\textbf{原地更新}，避免无意义的新 tensor 分配。

\subsubsection{组件：\texttt{AdamW}}
\paragraph{1. 作用、输入与输出}
实现带 decoupled weight decay 的 AdamW 优化器，输入是一组模型参数与超参数 \((lr,\beta_1,\beta_2,\epsilon,\lambda)\)。核心输出是一次 \texttt{step()} 后的参数更新。

\paragraph{2. 代码实现摘要}
\begin{lstlisting}
exp_avg.mul_(beta1).add_(grad, alpha=1.0 - beta1)
exp_avg_sq.mul_(beta2).addcmul_(grad, grad, value=1.0 - beta2)
bias_correction1 = 1.0 - beta1**step
bias_correction2 = 1.0 - beta2**step
p.data.mul_(1.0 - lr * weight_decay)
p.data.addcdiv_(corrected_exp_avg, denom, value=-lr)
\end{lstlisting}

\paragraph{3. 数学原理}
AdamW 的状态更新为
\[
m_t = \beta_1 m_{t-1} + (1-\beta_1) g_t,
\qquad
v_t = \beta_2 v_{t-1} + (1-\beta_2) g_t^2.
\]
偏置修正为
\[
\hat{m}_t = \frac{m_t}{1-\beta_1^t},
\qquad
\hat{v}_t = \frac{v_t}{1-\beta_2^t}.
\]
最终参数更新为
\[
\theta \leftarrow \theta (1 - lr \cdot \lambda)
 - lr \cdot \frac{\hat{m}_t}{\sqrt{\hat{v}_t} + \epsilon}.
\]
其中第一项就是 \textbf{decoupled weight decay}。

\paragraph{4. 工程技巧}
实现复用了 PyTorch \texttt{Optimizer} 基类的 \texttt{param\_groups} 与 \texttt{state} 机制，使接口与官方优化器一致。

\subsubsection{组件：\texttt{get\_lr\_cosine\_schedule}}
\paragraph{1. 作用、输入与输出}
根据当前 iteration \texttt{it} 返回学习率。输入包括最大/最小学习率、warmup 步数和 cosine 衰减总长度，输出是当前步的标量学习率。

\paragraph{2. 代码实现摘要}
\begin{lstlisting}
if warmup_iters > 0 and it < warmup_iters:
    return max_learning_rate * it / warmup_iters
if it >= cosine_cycle_iters:
    return min_learning_rate
progress = (it - warmup_iters) / (cosine_cycle_iters - warmup_iters)
cosine_decay = 0.5 * (1.0 + math.cos(math.pi * progress))
\end{lstlisting}

\paragraph{3. 数学原理}
该 schedule 是分段函数：
\[
\alpha_t =
\begin{cases}
\alpha_{\max}\cdot \frac{t}{T_w}, & 0 \le t < T_w, \\
\alpha_{\min} + \frac{1+\cos(\pi \cdot p_t)}{2}(\alpha_{\max}-\alpha_{\min}), & T_w \le t < T_c, \\
\alpha_{\min}, & t \ge T_c,
\end{cases}
\]
其中
\[
p_t = \frac{t-T_w}{T_c-T_w}.
\]

\paragraph{4. 工程技巧}
这里最关键的并不是代码量，而是\textbf{边界定义与测试语义一致}：第 \(0\) 步学习率应为 \(0\)，而不是 \(\alpha_{\max}/T_w\)。

\subsubsection{组件：\texttt{save\_checkpoint} / \texttt{load\_checkpoint}}
\paragraph{1. 作用、输入与输出}
保存并恢复训练状态。输入包括模型、优化器与 iteration，输出要么写入文件，要么恢复状态并返回被恢复的 iteration。

\paragraph{2. 代码实现摘要}
\begin{lstlisting}
checkpoint = {
    "model_state_dict": model.state_dict(),
    "optimizer_state_dict": optimizer.state_dict(),
    "iteration": iteration,
}
torch.save(checkpoint, out)
\end{lstlisting}
加载时：
\begin{lstlisting}
checkpoint = torch.load(src, map_location="cpu")
model.load_state_dict(checkpoint["model_state_dict"])
optimizer.load_state_dict(checkpoint["optimizer_state_dict"])
return checkpoint["iteration"]
\end{lstlisting}

\paragraph{4. 工程技巧}
这里采用 PyTorch 标准的 \texttt{state\_dict} 协议，优点是和模型/优化器内部结构解耦。

\subsection{注意力、位置编码与模型组装}

\subsubsection{组件：\texttt{scaled\_dot\_product\_attention}}
\paragraph{1. 作用、输入与输出}
给定
\[
Q \in \mathbb{R}^{\cdots \times n_q \times d_k},\quad
K \in \mathbb{R}^{\cdots \times n_k \times d_k},\quad
V \in \mathbb{R}^{\cdots \times n_k \times d_v},
\]
输出
\[
O \in \mathbb{R}^{\cdots \times n_q \times d_v}.
\]
\texttt{mask=True} 表示允许注意力流动，\texttt{False} 表示该位置必须被屏蔽。

\paragraph{2. 代码实现摘要}
\begin{lstlisting}
scores = torch.matmul(Q, K.transpose(-1, -2)) / math.sqrt(d_k)
if mask is not None:
    scores = scores.masked_fill(~mask, torch.finfo(scores.dtype).min)
attention_probs = softmax(scores, dim=-1)
return torch.matmul(attention_probs, V)
\end{lstlisting}

\paragraph{3. 数学原理}
\[
\mathrm{Attention}(Q,K,V)=\mathrm{softmax}\left(\frac{QK^\top}{\sqrt{d_k}}\right)V.
\]
mask 的实现方式是在 softmax 前把非法位置加上 \(-\infty\)，即：
\[
S_{ij} =
\begin{cases}
\frac{q_i^\top k_j}{\sqrt{d_k}}, & M_{ij}=\texttt{True}, \\
-\infty, & M_{ij}=\texttt{False}.
\end{cases}
\]

\paragraph{4. 工程技巧}
注意力实现支持任意数量的\textbf{batch-like 维度}，例如 \((batch, heads, seq, d)\)。

\subsubsection{组件：\texttt{rope}}
\paragraph{1. 作用、输入与输出}
RoPE 对 Query 或 Key 的最后一个维度做位置相关旋转，输入输出 shape 相同。

\paragraph{2. 代码实现摘要}
\begin{lstlisting}
x_even = in_query_or_key[..., ::2]
x_odd = in_query_or_key[..., 1::2]
rotated_even = x_even * cos - x_odd * sin
rotated_odd = x_even * sin + x_odd * cos
\end{lstlisting}

\paragraph{3. 数学原理}
把维度两两分组，记第 \(m\) 对维度对应角度为
\[
\phi_m(p)=\frac{p}{\Theta^{2m/d_k}},
\]
则对位置 \(p\) 上的向量分量 \((x_{2m},x_{2m+1})\) 做二维旋转：
\[
\begin{pmatrix}
x'_{2m} \\
x'_{2m+1}
\end{pmatrix}
=
\begin{pmatrix}
\cos \phi_m(p) & -\sin \phi_m(p) \\
\sin \phi_m(p) & \cos \phi_m(p)
\end{pmatrix}
\begin{pmatrix}
x_{2m} \\
x_{2m+1}
\end{pmatrix}.
\]

\paragraph{4. 工程技巧}
在多头注意力中，\textbf{head 维度被当作 batch-like 维度处理}，因此同一套旋转公式可以直接广播到所有 head。

\subsubsection{组件：\texttt{multihead\_self\_attention}}
\paragraph{1. 作用、输入与输出}
实现\textbf{因果多头自注意力}。输入为
\[
X \in \mathbb{R}^{\cdots \times T \times d_{\text{model}}},
\]
输出同样是
\[
\mathbb{R}^{\cdots \times T \times d_{\text{model}}}.
\]

\paragraph{2. 代码实现摘要}
\begin{lstlisting}
q = linear(..., weights=q_proj_weight, in_features=in_features)
k = linear(..., weights=k_proj_weight, in_features=in_features)
v = linear(..., weights=v_proj_weight, in_features=in_features)
q = q.reshape(..., num_heads, q_head_dim).transpose(-3, -2)
...
causal_mask = torch.tril(torch.ones(seq_len, seq_len, dtype=torch.bool, device=in_features.device))
attention_output = scaled_dot_product_attention(Q=q, K=k, V=v, mask=causal_mask)
\end{lstlisting}

\paragraph{3. 数学原理}
\[
\mathrm{MHA}(x)=W_O \cdot \mathrm{Concat}(\mathrm{head}_1,\dots,\mathrm{head}_h),
\]
其中每个 head 是
\[
\mathrm{head}_i = \mathrm{Attention}(Q_i, K_i, V_i).
\]
因果约束要求 query 位置 \(i\) 只能看见 \(j \le i\) 的 key，因此 mask 是下三角矩阵。

\paragraph{4. 工程技巧}
Q/K/V 的投影被实现为\textbf{三个 batched 线性变换}，而不是按 head 循环；这能更好利用张量并行。

\subsubsection{组件：\texttt{multihead\_self\_attention\_with\_rope}}
\paragraph{1. 作用、输入与输出}
它是普通因果 MHA 的 RoPE 扩展版本。与上一个组件不同的是：\textbf{只对 Q 和 K 应用 RoPE，不对 V 应用。}

\paragraph{2. 代码实现摘要}
\begin{lstlisting}
q = rope(..., in_query_or_key=q, token_positions=token_positions)
k = rope(..., in_query_or_key=k, token_positions=token_positions)
attention_output = scaled_dot_product_attention(Q=q, K=k, V=v, mask=causal_mask)
\end{lstlisting}

\paragraph{3. 数学原理}
\[
Q' = \mathrm{RoPE}(Q), \qquad K' = \mathrm{RoPE}(K), \qquad
\mathrm{Attention}(Q',K',V).
\]
这会让注意力分数显式携带相对位置信息。

\subsubsection{组件：\texttt{transformer\_block}}
\paragraph{1. 作用、输入与输出}
实现一个\textbf{pre-norm Transformer block}。输入输出 shape 都是
\[
\mathbb{R}^{batch \times seq \times d_{\text{model}}}.
\]

\paragraph{2. 代码实现摘要}
\begin{lstlisting}
x_norm = rmsnorm(..., weights=weights["ln1.weight"], in_features=in_features)
attn_out = multihead_self_attention_with_rope(..., in_features=x_norm, ...)
residual = in_features + attn_out
residual_norm = rmsnorm(..., weights=weights["ln2.weight"], in_features=residual)
ff_out = swiglu(..., in_features=residual_norm)
return residual + ff_out
\end{lstlisting}

\paragraph{3. 数学原理}
令输入为 \(x\)，则 block 的两个子层分别是：
\[
y = x + \mathrm{MHA}(\mathrm{RMSNorm}(x)),
\]
\[
z = y + \mathrm{FFN}(\mathrm{RMSNorm}(y)).
\]
这里 FFN 实际上就是前面定义的 SwiGLU。

\paragraph{4. 工程技巧}
采用 \textbf{pre-norm} 结构，而不是 post-norm。这样梯度路径更稳定，也更符合 handout 的要求。

\subsubsection{组件：\texttt{transformer\_lm}}
\paragraph{1. 作用、输入与输出}
这是最终的语言模型前向传播。输入是 token id 张量
\[
I \in \mathbb{N}^{B \times T},
\]
输出是 logits
\[
Z \in \mathbb{R}^{B \times T \times V}.
\]

\paragraph{2. 代码实现摘要}
\begin{lstlisting}
hidden = embedding(..., weights=weights["token_embeddings.weight"], token_ids=in_indices)
for layer_idx in range(num_layers):
    layer_weights = {...}
    hidden = transformer_block(..., weights=layer_weights, in_features=hidden)
hidden = rmsnorm(..., weights=weights["ln_final.weight"], in_features=hidden)
return linear(..., weights=weights["lm_head.weight"], in_features=hidden)
\end{lstlisting}

\paragraph{3. 数学原理}
定义
\[
H_0 = \mathrm{Embed}(I),
\]
第 \(\ell\) 层 block 为 \(B_\ell\)，则
\[
H_{\ell+1} = B_\ell(H_\ell), \qquad \ell=0,\dots,L-1.
\]
最后输出
\[
Z = H_L^{(\text{final norm})} W_{\text{lm}}^\top.
\]

\newpage
\section{测试过程、定位过程与最终结果}

\subsection{测试总流程}
测试策略采用\textbf{由局部到整体}的方式：
\begin{enumerate}[leftmargin=2em]
  \item 先通过 \texttt{py\_compile} 排除语法问题；
  \item 然后按组件粒度运行单测，例如 \texttt{test\_silu}、\texttt{test\_rmsnorm}、\texttt{test\_rope}；
  \item 对关键性能约束单独验证，如 \texttt{test\_train\_bpe\_speed}；
  \item 最后运行完整测试集 \texttt{pytest -q} 做集成回归。
\end{enumerate}

\subsection{关键问题的定位与修复}

\subsubsection{案例一：学习率调度 warmup 公式与测试定义不一致}
\paragraph{现象}
\texttt{test\_get\_lr\_cosine\_schedule} 初次失败，失败用例显示期望序列的开头是
\[
[0,\ 0.142857,\ 0.285714,\ \dots]
\]
而初版实现把 warmup 写成了
\[
\alpha_t = \alpha_{\max}\cdot\frac{t+1}{T_w},
\]
因此第 \(0\) 步学习率不是 \(0\)，而是 \(\alpha_{\max}/T_w\)。

\paragraph{定位}
通过直接阅读 \texttt{tests/test\_optimizer.py} 中的 \texttt{expected\_lrs}，可以明确看出测试采用的是
\[
\alpha_t = \alpha_{\max}\cdot\frac{t}{T_w}.
\]
因此问题不是余弦退火部分，而是 warmup 起点定义。

\paragraph{修复}
将 warmup 段修改为：
\begin{lstlisting}
if warmup_iters > 0 and it < warmup_iters:
    return max_learning_rate * it / warmup_iters
\end{lstlisting}
修复后，单测
\texttt{test\_get\_lr\_cosine\_schedule} 通过。

\subsubsection{案例二：BPE 训练性能瓶颈}
\paragraph{现象}
\texttt{test\_train\_bpe\_speed} 初次失败，测试内部记录到的训练时间约为 \(2.17\) 秒，而阈值是 \(1.5\) 秒。

\paragraph{定位}
初版 \texttt{train\_bpe} 每一轮都执行：
\begin{enumerate}[leftmargin=2em]
  \item 对全部 token sequence 重新统计 pair 频次；
  \item 对全部 token sequence 重新执行 merge。
\end{enumerate}
这使得复杂度近似变成“merge 轮数 \(\times\) 全语料重扫”，在小语料上也会明显超时。

\paragraph{修复}
最终版本改成了\textbf{增量统计}：
\begin{itemize}[leftmargin=2em]
  \item 用 \texttt{pair\_counts} 保存全局 pair 频次；
  \item 用 \texttt{pair\_to\_sequences} 保存 pair 到 sequence 的反向索引；
  \item 每轮只更新包含 \texttt{best\_pair} 的受影响序列。
\end{itemize}
修复后，\texttt{test\_train\_bpe\_speed} 通过。单独运行该测试命令时，pytest 报告时间为 \(0.31\) 秒；包含解释器与 pytest 启动开销的 wall-clock 时间约为 \(1.94\) 秒，但测试内部计时已经低于 \(1.5\) 秒阈值。

\subsubsection{案例三：GPT-2 参考对齐测试的可复现化}
\paragraph{现象}
tokenizer 对齐测试依赖 \texttt{tiktoken.get\_encoding("gpt2")}。为了让对齐测试在项目内可复现，需要保证 GPT-2 参考资源 \texttt{encoder.json} 与 \texttt{vocab.bpe} 可以被本地稳定读取。

\paragraph{处理方式}
最终方案是在 \texttt{tests/adapters.py} 顶部初始化本地 \texttt{TIKTOKEN\_CACHE\_DIR}，并把项目内已有的 GPT-2 fixture 写入缓存。其中：
\begin{itemize}[leftmargin=2em]
  \item \texttt{encoder.json} 直接使用 \texttt{tests/fixtures/gpt2\_vocab.json}
  \item \texttt{vocab.bpe} 使用 \texttt{tests/fixtures/gpt2\_merges.txt}，并补写 \texttt{\#version: 0.2} 头
\end{itemize}
这样对齐测试比较的是\textbf{实现行为}，而不是外部资源可达性。

\subsection{最终测试矩阵}
表 \ref{tab:all-tests} 总结了最终测试结果。所有测试均通过，只有 \texttt{test\_encode\_memory\_usage} 是测试中显式标记的 \texttt{xfail}，属于预期行为。

\footnotesize
\begin{longtable}{p{0.43\textwidth} p{0.12\textwidth} p{0.34\textwidth}}
\caption{最终测试矩阵}\label{tab:all-tests}\\
\toprule
测试名 & 结果 & 说明 \\
\midrule
\endfirsthead
\toprule
测试名 & 结果 & 说明 \\
\midrule
\endhead
\texttt{test\_get\_batch} & PASS & 数据采样、右移标签与 device 行为正确。 \\
\texttt{test\_linear} & PASS & 线性层矩阵乘法与 shape 对齐。 \\
\texttt{test\_embedding} & PASS & Embedding 查表正确。 \\
\texttt{test\_swiglu} & PASS & 前馈门控结构与参考权重对齐。 \\
\texttt{test\_scaled\_dot\_product\_attention} & PASS & 2D attention 正确。 \\
\texttt{test\_4d\_scaled\_dot\_product\_attention} & PASS & batch-like 维度广播正确。 \\
\texttt{test\_multihead\_self\_attention} & PASS & 因果 MHA 正确。 \\
\texttt{test\_multihead\_self\_attention\_with\_rope} & PASS & 带 RoPE 的 MHA 正确。 \\
\texttt{test\_transformer\_lm} & PASS & 完整 LM 前向通过。 \\
\texttt{test\_transformer\_lm\_truncated\_input} & PASS & 截断输入下的 LM 行为正确。 \\
\texttt{test\_transformer\_block} & PASS & pre-norm Block 正确。 \\
\texttt{test\_rmsnorm} & PASS & RMSNorm 数值对齐。 \\
\texttt{test\_rope} & PASS & 位置旋转编码正确。 \\
\texttt{test\_silu\_matches\_pytorch} & PASS & SiLU 与 PyTorch 对齐。 \\
\texttt{test\_softmax\_matches\_pytorch} & PASS & softmax 数值稳定且正确。 \\
\texttt{test\_cross\_entropy} & PASS & 交叉熵与参考对齐。 \\
\texttt{test\_gradient\_clipping} & PASS & global norm clipping 行为正确。 \\
\texttt{test\_adamw} & PASS & AdamW 与参考/官方实现对齐。 \\
\texttt{test\_get\_lr\_cosine\_schedule} & PASS & warmup + cosine 衰减正确。 \\
\texttt{test\_checkpointing} & PASS & checkpoint 存取一致。 \\
\texttt{test\_roundtrip\_empty} & PASS & 空字符串 roundtrip 正确。 \\
\texttt{test\_empty\_matches\_tiktoken} & PASS & 空输入与 GPT-2 参考一致。 \\
\texttt{test\_roundtrip\_single\_character} & PASS & 单字符 roundtrip 正确。 \\
\texttt{test\_single\_character\_matches\_tiktoken} & PASS & 单字符与参考一致。 \\
\texttt{test\_roundtrip\_single\_unicode\_character} & PASS & 单 Unicode 字符 roundtrip 正确。 \\
\texttt{test\_single\_unicode\_character\_matches\_tiktoken} & PASS & 单 Unicode 字符与参考一致。 \\
\texttt{test\_roundtrip\_ascii\_string} & PASS & ASCII 串 roundtrip 正确。 \\
\texttt{test\_ascii\_string\_matches\_tiktoken} & PASS & ASCII 串与参考一致。 \\
\texttt{test\_roundtrip\_unicode\_string} & PASS & Unicode 串 roundtrip 正确。 \\
\texttt{test\_unicode\_string\_matches\_tiktoken} & PASS & Unicode 串与参考一致。 \\
\texttt{test\_roundtrip\_unicode\_string\_with\_special\_tokens} & PASS & 特殊 token 混合输入 roundtrip 正确。 \\
\texttt{test\_unicode\_string\_with\_special\_tokens\_matches\_tiktoken} & PASS & 与参考一致。 \\
\texttt{test\_overlapping\_special\_tokens} & PASS & 重叠 special token 匹配正确。 \\
\texttt{test\_address\_roundtrip} & PASS & 地址文本 roundtrip 正确。 \\
\texttt{test\_address\_matches\_tiktoken} & PASS & 地址文本与参考一致。 \\
\texttt{test\_german\_roundtrip} & PASS & 德文文本 roundtrip 正确。 \\
\texttt{test\_german\_matches\_tiktoken} & PASS & 德文文本与参考一致。 \\
\texttt{test\_tinystories\_sample\_roundtrip} & PASS & TinyStories 样本 roundtrip 正确。 \\
\texttt{test\_tinystories\_matches\_tiktoken} & PASS & TinyStories 样本与参考一致。 \\
\texttt{test\_encode\_special\_token\_trailing\_newlines} & PASS & 特殊 token 邻接换行的行为正确。 \\
\texttt{test\_encode\_special\_token\_double\_newline\_non\_whitespace} & PASS & 特殊 token 后双换行且后接文本时行为正确。 \\
\texttt{test\_encode\_iterable\_tinystories\_sample\_roundtrip} & PASS & 流式编码 roundtrip 正确。 \\
\texttt{test\_encode\_iterable\_tinystories\_matches\_tiktoken} & PASS & 流式编码与参考一致。 \\
\texttt{test\_encode\_iterable\_memory\_usage} & PASS & 流式编码满足内存约束。 \\
\texttt{test\_encode\_memory\_usage} & XFAIL & 预期失败：\texttt{encode} 不是内存优化接口。 \\
\texttt{test\_train\_bpe\_speed} & PASS & BPE 训练速度满足阈值。 \\
\texttt{test\_train\_bpe} & PASS & BPE merges / vocab 正确。 \\
\texttt{test\_train\_bpe\_special\_tokens} & PASS & 特殊 token 边界处理正确。 \\
\bottomrule
\end{longtable}
\normalsize

\subsection{最终汇总}
最终完整测试命令
\begin{lstlisting}
./.venv/bin/python -m pytest -q
\end{lstlisting}
得到的最终结果是
\begin{lstlisting}
47 passed, 1 xfailed in 13.04s
\end{lstlisting}
包含解释器与 pytest 启动开销的 wall-clock 时间约为
\[
15.04\ \text{s}.
\]

\newpage
\section{关于 LLM 实现的收获与总结}

\subsection{最重要的技术收获}
\begin{enumerate}[leftmargin=2em]
  \item \textbf{LLM 并不是一个单独的“模型文件”，而是完整的系统链路。} \\
  从 tokenizer、batch 构造、损失函数、优化器到注意力与 Block 组装，任何一层出错都可能让整体训练失败。

  \item \textbf{工程正确性与数学正确性同等重要。} \\
  例如 \texttt{softmax} 与 \texttt{cross\_entropy} 的数学公式都很短，但如果不做数值稳定处理，实际训练中会立刻出问题。

  \item \textbf{性能问题不能只靠“更快的机器”掩盖。} \\
  \texttt{train\_bpe} 的主要瓶颈来自算法结构，而不是单条指令慢。把“全量重扫”改成“增量维护索引”以后，性能才真正改善。

  \item \textbf{自注意力的核心难点在 shape 管理，而不是公式本身。} \\
  公式
  \[
  \mathrm{softmax}\left(\frac{QK^\top}{\sqrt{d_k}}\right)V
  \]
  并不复杂，真正容易错的是 head 拆分、batch-like 维度广播、causal mask 的形状，以及 RoPE 应该施加在 \(Q/K\) 而不是 \(V\) 上。

  \item \textbf{Transformer 的稳定训练依赖若干“局部正确的组件约定”。} \\
  比如 pre-norm、global gradient clipping、AdamW 的 decoupled weight decay、cosine LR schedule 的 warmup 边界，都是看似局部但对整体训练质量有直接影响的设计。
\end{enumerate}

\subsection{关于 LLM 的理解深化}
通过手写这一版 Assignment 1，我对 LLM 的理解从“知道结构”进一步推进到了“知道数据和张量是怎样真正流动的”。例如，过去在阅读论文时，\texttt{RoPE}、\texttt{MHA}、\texttt{RMSNorm} 往往被视作已经打包好的模块；但在这次实现中，必须亲自处理如下事实：
\begin{itemize}[leftmargin=2em]
  \item tokenization 实际上决定了后续模型要学习的离散表示空间；
  \item attention 的 mask 语义、softmax 的数值稳定性，会直接影响训练是否可靠；
  \item Block 与 LM 的实现不是“搭积木”那么简单，而是要求每一层 shape 都严格一致；
  \item 训练系统的可恢复性（checkpoint）、可控性（lr schedule）、稳定性（gradient clipping）和模型前向本身同样重要。
\end{itemize}

\subsection{总结}
从结果上看，这次 Assignment 1 完成了一个\textbf{从 tokenizer 到 Transformer LM 前向、从 batch 到 optimizer step 的完整最小闭环}。如果只从代码规模看，这些函数都不长；但真正困难的部分在于：
\begin{enumerate}[leftmargin=2em]
  \item 把 handout 中的数学定义精确地落实为张量操作；
  \item 保证所有中间接口（shape、dtype、mask 语义、special token 边界）完全一致；
  \item 用测试驱动的方式把 correctness、numerical stability 和 performance 都收口。
\end{enumerate}

因此，这份作业最有价值的地方，并不只是“写出了一个小 Transformer”，而是把\textbf{LLM 系统实现中的关键思想、关键接口和关键调试方法}都走了一遍。对后续继续做训练脚本、解码、实验日志和更大规模实验来说，这一版代码已经形成了一个可靠的基础。

\end{document}
